\clearpage
\phantomsection
\addcontentsline{toc}{chapter}{Приложение А. Воспроизведение измерений}
\chapter*{Приложение А}
\vspace{-12pt}
{\centering\singlespacing\textbf{Воспроизведение измерений}\par}
\vspace{12pt}

Все результаты отчёта хранятся в репозитории~\cite{mz77-repo} в каталоге
\texttt{bench/log}: подкаталог на каждый компьютер и шаг, в нём исходные
журналы (\texttt{raw/}) и файл \texttt{EXPLANATION.md} с описанием условий
измерения. Указатель на все результаты "--- файл
\texttt{doc/optimisations/RESULTS.md}, описание каждой оптимизации "---
файлы \texttt{doc/optimisations/*.md}.

Для измерений нужны Linux на x86-64, GCC, autoconf, automake, libtool,
gnulib, файлы разработки GTK~4, \texttt{perf} и \texttt{taskset}. Счётчики
процессора должны быть доступны пользователю:
\texttt{sysctl kernel.perf\_event\_paranoid=1}. Ноутбук должен быть
подключён к сети, компьютер не должен быть нагружен.

Сборка и загрузка корпуса:
\begin{lstlisting}[language=bash,numbers=none]
./bootstrap
./configure --enable-release
make -j$(nproc)
make corpus
\end{lstlisting}

Измерение всех шагов в тактах (таблица~\ref{tab:final-cycles}): сценарий
собирает каждый шаг из его коммита и пишет \texttt{steps.txt} в
\texttt{bench/results/КОММИТ/}.
\begin{lstlisting}[language=bash,numbers=none]
ROUNDS=5 CPU=3 bench/steps.sh
\end{lstlisting}

Полное измерение одного шага (время, профиль, контрольные суммы архивов):
\begin{lstlisting}[language=bash,numbers=none]
git checkout 5d53dc4
./configure --enable-release && make clean && make -j$(nproc)
make pgo
MACHINE=mybox bench/campaign.sh 07-lto-pgo
\end{lstlisting}

Отдельные инструменты: \texttt{make bench} "--- время на всём корпусе,
\texttt{make bench-compare} "--- сборки с разными флагами компилятора,
\texttt{make profile} "--- профиль одного файла (gprof, perf, Cachegrind,
Callgrind). Переменная \texttt{MZ77\_MATCH=word} или
\texttt{MZ77\_MATCH=avx2} выбирает реализацию сравнения.

Данные для графиков этого отчёта извлекаются из \texttt{bench/log}
сценарием \texttt{archive/report-2026-10/data/extract.py}; измерение
зависимости от размера входа повторяет сценарий
\texttt{archive/report-2026-10/data/scaling.sh}.
